% Source: https://artofproblemsolving.com/community/c3h1565293_interesting_problem_relating_to_sequences
% Stable identity: aops-topic-1565293-q01
\textbf{What $a_n$ counts.} The number $n-2\left\lfloor\dfrac {n}{2}\right\rfloor$ is the remainder of $n$ on division by $2$, i.e.\ the last binary digit of $n$, and $\left\lfloor\dfrac {n}{2}\right\rfloor$ is $n$ with its last binary digit deleted. So the recursion adds the last binary digit and moves on to the number formed by the remaining digits. By induction on $n$ (with $a_0=0$), $a_n$ is the number of $1$s in the binary representation of $n$.

\textbf{Binary form of $N$.} Since $2^{24}-1=\left(2^3\right)^8-1=(2^3-1)\left(2^{21}+2^{18}+\cdots+2^3+1\right)=7\sum_{j=0}^{7}8^j$,
\[
N=\dfrac{5}{7}\left(2^{24}-1\right)=5\sum_{j=0}^{7}8^{j}=\sum_{j=0}^{7}101_2\cdot2^{3j}.
\]
Because $5=101_2$ has only three binary digits, the shifted copies do not overlap, so
\[
N=\underbrace{101\,101\,101\cdots101_2}_{8\text{ blocks}} .
\]
Each block contributes two $1$s, so $a_N=8\cdot2=\boxed{16}$.
